\documentclass[12pt,a4paper]{article}
\usepackage[utf8]{inputenc}
\usepackage[ngerman]{babel}
\usepackage[margin=1.5cm]{geometry}
\usepackage{multicol}
\usepackage{parskip}
\usepackage{enumitem}
\usepackage{amsmath}

\title{Spickzettel Biologie -- Skelett \& Knochen}
\author{7. Klasse Realschule, Baden-Württemberg}
\date{}

\begin{document}
	
	\maketitle
	\begin{multicols}{2}
		
		\section*{Aufbau und Funktion des Skeletts}
		Das menschliche Skelett besteht beim Erwachsenen aus etwa \textbf{206 Knochen}. Es bildet das stabile, aber leichte und flexible Gerüst des Körpers.
		
		\subsection*{Die vier Hauptfunktionen des Skeletts}
		\begin{enumerate}[noitemsep]
			\item \textbf{Stützfunktion:} Gibt dem Körper Form, Größe und Halt. 
			\item \textbf{Schutzfunktion:} Schützt innere Organe vor Verletzungen.
			\begin{itemize}[noitemsep]
				\item Schädel $\rightarrow$ schützt das Gehirn
				\item Brustkorb (Rippen + Brustbein) $\rightarrow$ schützt Herz und Lunge
				\item Becken $\rightarrow$ schützt Teile des Darms, Harnblase und Gebärmutter
				\item Wirbelsäule $\rightarrow$ schützt das Rückenmark
			\end{itemize}
			\item \textbf{Bewegungsfunktion:} Zusammen mit Muskeln, Sehnen und Bändern ermöglicht das Skelett Bewegungen.
			\item \textbf{Blutbildung und Speicherfunktion:} Im Knochenmark werden Blutzellen gebildet; Knochen speichern Kalzium und Phosphat.
		\end{enumerate}
		
		\subsection*{Einteilung des Skeletts}
		Das Skelett wird in zwei Hauptgruppen unterteilt:
		\begin{itemize}[noitemsep]
			\item \textbf{Achsenskelett:} Schädel, Wirbelsäule, Brustkorb (Rippen, Brustbein)
			\item \textbf{Extremitätenskelett (Gliedmaßenskelett):} Schultergürtel, Arme, Beckengürtel, Beine
		\end{itemize}
		
		\section*{Die Wirbelsäule}
		Die Wirbelsäule besteht aus \textbf{7 Halswirbeln, 12 Brustwirbeln, 5 Lendenwirbeln, Kreuzbein und Steißbein}. 
		Zwischen den Wirbeln liegen die \textbf{Bandscheiben}. Sie wirken wie Stoßdämpfer und federn Stöße ab (z.B. beim Laufen oder Springen).
		
		\section*{Aufbau eines Knochens (am Beispiel Röhrenknochen)}
		Ein Röhrenknochen (z.B. Oberarm- oder Oberschenkelknochen) besteht aus:
		\begin{itemize}[noitemsep]
			\item \textbf{Diaphyse (Schaft):} Der lange, mittlere Teil des Knochens.
			\item \textbf{Epiphysen (Gelenkenden):} Die beiden verdickten Enden.
			\item \textbf{Wachstumsfugen:} Zwischen Schaft und Gelenkenden; hier findet das Längenwachstum statt. Sind sie verknöchert, ist das Wachstum abgeschlossen.
			\item \textbf{Knochenmarkhöhle:} Im Inneren des Schaftes; enthält das blutbildende Knochenmark.
		\end{itemize}
		
		\subsection*{Knochenarten}
		Man unterscheidet verschiedene Knochenformen:
		\begin{itemize}[noitemsep]
			\item \textbf{Röhrenknochen (lange Knochen):} Wirken als Hebel, z.B. Oberarm-, Oberschenkelknochen, Schienbein.
			\item \textbf{Platte Knochen:} Z.B. Schädelknochen, Schulterblatt.
			\item \textbf{Kurze Knochen:} Z.B. Hand- und Fußwurzelknochen.
		\end{itemize}
		
		\subsection*{Knochenzusammensetzung}
		Knochen bestehen aus einer gummiartigen Grundmasse (Elastin/Kollagen), in die harter \textbf{Knochenkalk} (Kalziumsalze) eingelagert ist. Diese Mischung macht sie \textbf{stabil und gleichzeitig elastisch}.
		
		\section*{Gelenke}
		Gelenke verbinden mindestens zwei Knochen miteinander und ermöglichen Bewegungen. Man unterscheidet:
		\begin{itemize}[noitemsep]
			\item \textbf{Bewegliche Gelenke:} Z.B. Schulter-, Hüft-, Kniegelenk.
			\item \textbf{Unbewegliche Gelenke:} Z.B. Schädel- oder Beckennähte.
		\end{itemize}
		\vspace{0.2cm}
		\noindent \textbf{Wichtige Gelenktypen:}
		\begin{itemize}[noitemsep]
			\item \textbf{Kugelgelenk:} Größte Beweglichkeit (nach allen Seiten), z.B. Schulter- und Hüftgelenk.
			\item \textbf{Scharniergelenk:} Bewegungen nur in eine Richtung (beugen/strecken), z.B. Ellenbogen- und Kniegelenk.
		\end{itemize}
		
		\section*{Wichtige Knochen -- Merkwissen}
		\begin{itemize}[noitemsep]
			\item \textbf{Längster Knochen:} Oberschenkelknochen (ca. 45--50 cm).
			\item \textbf{Kleinste Knochen:} Gehörknöchelchen im Mittelohr (Steigbügel, Amboss, Hammer).
			\item \textbf{Anzahl der Knochen bei Babys:} Ca. 300, da viele Knochen erst später zusammenwachsen.
			\item \textbf{Gewichtsanteil:} Knochen machen ca. 14\,\% des Körpergewichts aus.
			\item \textbf{Wichtige Knochen erkennen:} Schädel, Wirbelsäule, Brustbein, Rippen, Schlüsselbein, Schulterblatt, Oberarmknochen, Elle, Speiche, Becken, Oberschenkelknochen, Kniescheibe, Schienbein, Wadenbein.
		\end{itemize}
		
		\section*{Tipps für den Test}
		\begin{enumerate}[noitemsep]
			\item Lerne die \textbf{Funktionen} des Skeletts auswendig (Stütze, Schutz, Bewegung, Blutbildung/Speicher).
			\item Übe, die wichtigsten Knochen an einer Skelettzeichnung zu \textbf{beschriften} (siehe Merkliste oben).
			\item Verstehe den \textbf{Aufbau eines Röhrenknochens} (Schaft, Gelenkenden, Wachstumsfugen).
			\item Kenne die \textbf{Einteilung} in Achsen- und Extremitätenskelett.
			\item Wisse, welche Knochen \textbf{welche Organe schützen}.
		\end{enumerate}
		
	\end{multicols}
\end{document}